
\chapter{Thermodynamics}


\section{Properties}
Thermodynamics properties can categorized as either {\bf intensive} 
or {\bf extensive}.  
\begin{itemize}
\item Intensive properties are independent of, and therefore not 
proportional to, the size or sample of the material.
Intensive properties are also referred to as {\bf intrinsic} properties
or, equivalently, as {\bf bulk} properties.
\item In contrast, extensive properties are 
dependent on, and therefore proportional
to, the size or sample of the material.  
Extrinsic properties are also referred to as {\bf extrinsic} properties.
\end{itemize}

Figure~\ref{fig:intrinsic_props} demonstrates the independence of intrinsic
properties and the dependence of extrinsic properties on the control volume, 
designated with the dashed lines.  The original control volume, on the left
side of the figure, is divided into two equal-sized control volumes, as 
shown on the 
right side of the figure.  The intrinsic properties do not depend on the 
size of the control volume.  The extrinsic properties, however, are now
scaled by a factor of $\half$.
\begin{figure}[htb]
  \begin{center}
  %\begin{overpic}[width=\textwidth, grid, tics=10]{fig/intrinsic_props}
  \begin{overpic}[width=\textwidth,natwidth=540,natheight=180]{fig/intrinsic_props.pdf}
    \put( 0, 26){\bf{Intrinsic} (intensive, bulk)}
    \put( 0, 23){pressure, density, temperature}
    \put(48, 24){$P, \rho, T$}
    \put(74.5, 24){$P, \rho, T$}
    \put(88.5, 24){$P, \rho, T$}
    \put( 0, 17){\bf{Extrinsic} (extensive)}
    \put( 0, 14){mass, volume, energy}
    \put(47.5, 15.5){$m, V, E$}
    \put(72.5, 15.2){$\dst\frac{m}{2}, \frac{V}{2}, \frac{E}{2}$}
    \put(87., 15.2){$\dst\frac{m}{2}, \frac{V}{2}, \frac{E}{2}$}
    \put( 0,  8){\bf{Specific} (extrinsic/mass)}
    \put( 0,  5){specific volume, specific energy}
    \put(49.5,  7){$\svol, \senergy$}
    \put(76.3,  7){$\svol, \senergy$}
    \put(90.5,  7){$\svol, \senergy$}
  \end{overpic}
  \end{center}
  \caption{Thermodynamic properties categorized as intrinsic, extrinsic, and
  specific.}
  \label{fig:intrinsic_props}
\end{figure}

It is often convenient to speak of extrinsic properties in away that makes
them {\em independent} of size or sample of the material.  
To achieve this, the
extrinsic properties are divided by the mass of the sample. 
This ratio is referred to 
as the specific extrinsic property.  
Two common examples of specific properties
are the {\bf specific volume},
\be
\svol \defe \frac{V}{m} = \frac{1}{\rho},
\label{eq:specific_volume}
\ee
and
{\bf specific energy},
\be
\senergy \defe \frac{E}{m}.
\ee


\section{First Law of Thermodynamics}
\begin{itemize}
\item A closed system can exchange heat and work
with its surroundings, but cannot exchange mass.

\item An open system can exchange heat, work, and mass with its surroundings.
\end{itemize}

The first law for a closed system states the change in internal
energy $\Delta U$ is equal to the net energy $Q$ added to the system as heat
less the net work $W$ done {\bf by} the system.
\be
\Delta U = Q - W.
\ee
This was the sign convention used by Joule and Clausius.
A different sign convention, used by Planck, stated that 
the change in internal energy $\Delta U$ is equal to the net energy $Q$
added to the system as heat plus the net work done {\bf on} the system.
\be
\Delta U = Q + W.
\ee
Thus, all net transfers {\bf to} the system are {\bf positive}.  In 
contrast, all net transfers {\bf from} the system are {\bf negative}.

Then, the conservation of energy principle for open systems can be
written as
\be
\Delta E = \Delta (K + U) = Q + W,
\ee
where the total energy $E$ is the sum of the kinetic energy $K$ and 
the internal energy $U$.
The change in total energy is equal to the net energy $Q$ added to the system
as heat plus the net work $W$ done {\bf on} the system.

Next, we introduce
\begin{itemize}
\item $d(\cdot)$ as the {\bf exact differential}, which does not depend on the 
path taken through space of the thermodynamic parameters and depends only
on the initial states, and
\item $\delta(\cdot)$ as the {\bf inexact differential}, which depends on the
path taken through space of thermodynamic parameters.
\end{itemize}
The balance law then states
\be
dE = \delta Q + \delta W.
\ee

\subsubsection{First Law in Rate Form}
We will be interested in the change in energy over a unit of time.
Thus, it is useful to write the First Law of Thermodynamics in
terms of time derivatives.
\be
\frac{dE}{dt} = \delta \dot Q + \delta \dot W
\ee


\subsubsection{Continuum Description}
Now we write explicit continuum descriptions for each of the 
terms in the balance law
\be
\frac{d}{dt} ( K + U ) = \delta \dot Q + \delta \dot W.
\ee
The kinetic energy is given by
\be
K = \int_V \half \; \rho(\vx, t) \; \vv(\vx, t) \cdot \vv(\vx, t) \; dV 
\ee

\section{Second Law of Thermodynamics}
The Second Law states that the infinitesimal increment of entropy $dS$ 
of a system from an infinitesimal transfer of heat $\delta Q$, divided
by the common temperature of the system and the surrounding that supplies the
heat, satisfies
\be
dS \geq \dst\frac{\delta Q}{T}.
\ee
For a reversible process, 
\be
dS = \dst\frac{\delta Q}{T}.
\ee
For an irreversible process, 
\be
dS > \dst\frac{\delta Q}{T}.
\ee

\section{First and Second Laws Combined}
For a reversible process, the first and second laws can be combined as
\be
\boxed{dE = T dS + \delta W, \hspace{1.0cm}\mbox{\bf 
(1st and 2nd Law, Reversible Form)}.}
\ee

\section{Thermodynamic Potentials}
\subsection{Variables}

\begin{tabular}{cll}
$P$ & pressure &  (N/m$^2$) \\
$V$ & volume & (m$^3$)\\
$T$ & temperature & (K) \\
$S$ & entropy & (J/K) 
\end{tabular}

\subsection{Potentials}
Thermodynamic potentials consider a thermodynamically closed system, 
thus the energy $E = (K + U)$ is considered only as internal energy $U$.

\begin{tabular}{clll}
$U$ & & internal energy & (J) \\
$\helmholtz$ &$\defe \; U - TS$ & Helmholtz free energy & (J) \\
$\gibbs$ &$\defe \; U + PV - TS$ & Gibbs free energy & (J) \\
$\enthalpy$ &$\defe \; U + PV$ & enthalpy & (J) 
\end{tabular}

\vspace{2em}
The {\bf Helmholtz} free energy $\helmholtz$ is the Legendre transform
of the internal energy $U$, where the entropy $S$ is replaced by
the temperature $T$ as the independent variable.  

\subsection{Differential Forms of Potentials}
\begin{tabular}{rcrcrl}
$dU$ &=& $ T \;dS$ & $-$ & $P \;dV$   &  \\
$d\helmholtz$ &=& $-P \;dV$ & $-$ & $S \;\;dT$ &  \\
$d\gibbs$ &=& $-S \;dT$ & $+$ & $V \;dP$   &  \\
$d\enthalpy$ &=& $ V \;dP$ & $+$ & $T \;dS$ &  
\end{tabular}

\begin{figure}[htb]
  \begin{center}
  %\begin{overpic}[width=0.7\textwidth, grid, tics=10]{fig/thermo_potentials}
  \begin{overpic}[width=0.7\textwidth,natwidth=198,natheight=198]{fig/thermo_potentials.pdf}
    \put(48, 61){$dU$}
    \put(35, 48){$d\helmholtz$}
    \put(48, 35){$d\gibbs$}
    \put(60, 48){$d\enthalpy$}
    \put(30, 65){$-dV$}
    \put(14, 84){$P$}
    \put(30, 32){$-dT$}
    \put(14, 13){$S$}
    \put(64, 32){$dP$}
    \put(84, 13){$V$}
    \put(64, 65){$dS$}
    \put(84, 84){$T$}
    \put(30, 76.5){\footnotesize negative}
    \put(58, 76.5){\footnotesize positive}
  \end{overpic}
  \end{center}
  \caption{Mnemonic for the thermodynamic potentials
  in differential form.}
  \label{fig:thermo_potentials}
\end{figure}

\section{Maxwell Relations}
The Maxwell relations are a statement of equality between
{\em first} derivative of a function because of a quality possessed by 
the {\em second} derivative of that function.  Specifically, 
the Maxwell relations equate {\em first} derivatives of the thermodynamic
potential functions $U, \helmholtz, \gibbs,$ and $\enthalpy$, 
because the {\em second} derivates of those
same functions does not depend on the order of differentiation
with respect to the thermodynamic variables $P, V, T,$ and $S$.

The Maxwell relations make use of the {\bf Schwarz theorem},
see~(\ref{eq:schwarz}).
With this context in mind, consider now the {\em first} derivatives 
of the thermodynamic potentials with respect to their natural variables.
Specifically, we can write the four thermodynamic potentials as functions
of their natural variables:
\begin{align}
 U &= U(S, V), \\
 \helmholtz &= \helmholtz(V, T), \\
 \gibbs &= \gibbs(T, P), \\
 \enthalpy &= \enthalpy(P, S). 
\end{align}
Next, we take {\em first} derivatives of these potentials:
\begin{align}
\left( \frac{\partial U}{\partial S} \right)_V = T, & \hspace{1.5cm}
\left( \frac{\partial U}{\partial V} \right)_S = -P, \\ 
& \nonumber \\ 
%
\left( \frac{\partial \helmholtz}{\partial V} \right)_T = -P, & \hspace{1.5cm}
\left( \frac{\partial \helmholtz}{\partial T} \right)_V = -S, \\
& \nonumber \\
%
\left( \frac{\partial \gibbs}{\partial T} \right)_P = -S, & \hspace{1.5cm}
\left( \frac{\partial \gibbs}{\partial P} \right)_T = V, \\ 
& \nonumber \\ 
%
\left( \frac{\partial \enthalpy}{\partial P} \right)_S = V, & \hspace{1.5cm}
\left( \frac{\partial \enthalpy}{\partial S} \right)_P = T. 
\end{align}
Finally, we take {\em second} derivatives of the potential functions
with respect to the natural variables, which
allows for the {\em first} derivatives of the natural variables with 
respect to the natural variables themselves to be equated, as follows:
\be
 \frac{\partial^2 U}{\partial V \; \partial S} 
 \;\;\implies \;\;
 + \left(\frac{ \partial T}{\partial V}\right)_S 
 \;\; = \;\;
 - \left(\frac{ \partial P}{\partial S}\right)_V
\ee
\be
 \frac{\partial^2 \helmholtz}{\partial T \; \partial V} 
 \;\;\implies \;\;
 - \left(\frac{ \partial P}{\partial T}\right)_V 
 \;\; = \;\;
 - \left(\frac{ \partial S}{\partial V}\right)_T
\ee
\be
 \frac{\partial^2 \gibbs}{\partial P \; \partial T} 
 \;\;\implies \;\;
 - \left(\frac{ \partial S}{\partial P}\right)_T 
 \;\; = \;\;
 + \left(\frac{ \partial V}{\partial T}\right)_P
\ee
\be
 \frac{\partial^2 \enthalpy}{\partial S \; \partial P} 
 \;\;\implies \;\;
 + \left(\frac{ \partial V}{\partial S}\right)_P 
 \;\; = \;\;
 + \left(\frac{ \partial T}{\partial P}\right)_S
\ee
The foregoing equations are useful for eliminating entropy $S$, which 
is difficult to measure in a laboratory setting, in favor of the 
other natural variables pressure $P$, volume $V$, and temperature $T$, 
all of which are relatively easier to measure.


